%%%%%%%%%%%%%%%%%%%%%%%%%%%%%%%%%%%%%%%%%%%%%%%%%%%%%%%%%%%%%%%%%%%%%%
% Lớp 10 | Hàm số bậc hai và đồ thị | Hàm số và đồ thị
% Dạng 0D3?1-6: Tính chẵn, lẻ
%%%%%%%%%%%%%%%%%%%%%%%%%%%%%%%%%%%%%%%%%%%%%%%%%%%%%%%%%%%%%%%%%%%%%%


% nguồn: TuLuan_Lop10_Dot2 (nhập 14/11/2026)
\begin{bt}%[0D3H1-6]%[Huỳnh Đức Vũ]%[0D1B1]%[ma:8c426010]
	Xét tính chẵn lẻ của hàm số $f(x)=\dfrac{3x^2+5}{x^3}$.
	\loigiai {Tập xác định  $\mathscr{D}=\mathbb{R}\backslash \{0\}$.\\
		Với mọi $x\in \mathscr{D}$ ta có
		\begin{itemize}
			\item $-x \in \mathscr{D}$.
			\item $f(-x)=\dfrac{3(-x)^2+5}{(-x)^3}=-\dfrac{3x^2+5}{x^3}=-f(x)$.
		\end{itemize}
		Vậy hàm số đã cho là hàm số lẻ.}
\end{bt}

% nguồn: TuLuan_Lop10_Dot2 (nhập 14/11/2026)
\begin{bt}%[0D3H1-6]%[Huỳnh Đức Vũ]%[0D1B1]%[ma:13709b86]
	Xét tính chẵn lẻ của hàm số $f(x)=\dfrac{\sqrt{1-x}-\sqrt{1+x}}{|x-2022|+|x+2022|}$.
	\loigiai {Tập xác định  $\mathscr{D}=[-1;1]$.\\
		Với mọi $x\in \mathscr{D}$ ta có
		\begin{itemize}
			\item $-x \in \mathscr{D}$.
			\item $f(x)=\dfrac{\sqrt{1+x}-\sqrt{1-x}}{|-x-2022|+|-x+2022|}=-\dfrac{\sqrt{1-x}-\sqrt{1+x}}{|x+2022|+|x-2022|}=-f(x)$.
		\end{itemize}
		Vậy hàm số đã cho là hàm số lẻ.}
\end{bt}

% nguồn: TuLuan_Lop10_Dot2 (nhập 14/11/2026)
\begin{bt}%[0D3H1-6]%[Huỳnh Đức Vũ]%[0D1B1]%[ma:9fcb05dc]
	Xét tính chẵn lẻ của hàm số $f(x)=\dfrac{\sqrt{9-x^2}}{x^4-1}$.
	\loigiai {Tập xác định  $\mathscr{D}=[-3;3]\setminus\{-1;1\}$.\\
		Với mọi $x\in \mathscr{D}$ ta có
		\begin{itemize}
			\item $-x \in \mathscr{D}$.
			\item $f(-x)=\dfrac{\sqrt{9-(-x)^2}}{(-x)^4-1}=\dfrac{\sqrt{9-x^2}}{x^4-1}$.
		\end{itemize}
		Vậy hàm số đã cho là hàm số chẵn.}
\end{bt}

% nguồn: TuLuan_Lop10_Dot2 (nhập 14/11/2026)
\begin{bt}%[0D3H1-6]%[Huỳnh Đức Vũ]%[0D1B1]%[ma:5febd8d8]
	Xét tính chẵn lẻ của hàm số $f(x)=\dfrac{\sqrt{9+x^2}}{x^3-1}$.
	\loigiai {Tập xác định  $\mathscr{D}=\mathbb{R}\setminus\{1\}$.\\
		Vì $x_0=-1\in \mathscr{D}$ nhưng $-x_0=1 \notin \mathscr{D}$ nên tập xác định $\mathscr{D}$ không phải là tập đối xứng.\\
		Vậy hàm số đã cho không phải là hàm số chẵn, không phải là hàm số lẻ.}
\end{bt}

% nguồn: TuLuan_Lop10_Dot2 (nhập 14/11/2026)
\begin{bt}%[0D3H1-6]%[Huỳnh Đức Vũ]%[0D1B1]%[ma:621ddd37]
	Xét tính chẵn lẻ của hàm số $f(x)=\dfrac{3x+\sqrt{3+x^2}}{|x|+2}$.
	\loigiai {Tập xác định  $\mathbb{R}$.\\
		Vì $f(-1)=-\dfrac{1}{3}\ne \pm f(1)= \dfrac{5}{3}$ nên hàm số đã cho không phải là hàm số chẵn, không phải là hàm số lẻ.}
\end{bt}

% nguồn: TuLuan_Lop10_Dot2 (nhập 14/11/2026)
\begin{bt}%[0D3H1-6]%[Dự án đề kiểm tra HKI NH22-23 - Phạm Phương]%[THCS&THPT Lương Thế Vinh - Sở Hà Nội]%[ma:4f7374e8]
	Cho hàm số $y=f(x)=x^{2023}+ax^{2019}+3$ thỏa mãn $f(2022)=2021$. Tính $f(-2022)$.
	\loigiai{
	Ta có $f(2022)=2021 \Leftrightarrow 2022^{2023}+a\cdot2022^{2019}+3=2021 \Leftrightarrow 2022^{2023}+a\cdot2022^{2019}=2018$.
	Khi đó
	\allowdisplaybreaks{
		\begin{eqnarray*}
			f(-2022)&= & (-2022)^{2023}+a\cdot(-2022)^{2019}+3\\
			&= & -\left(2022^{2023}+a\cdot2022^{2019}\right)+3 \\
			&= & -2018+3=-2015.
	\end{eqnarray*} }
Vậy $f(-2022)=-2015$.
}
\end{bt}
